\section{Results}
\label{sec:results}

\noindent\textit{Every table in this section is a generated fragment
(\texttt{paper/tables/<name>.tex}, one \texttt{booktabs} body per hook)
produced by the analysis stage from the per-cell results files; every
figure is a generated \texttt{figures/<name>.pdf}. No number in this
section is typed by hand. The shift-regime and placebo columns are final:
each is read from its regime's two results files
(\texttt{results/<regime>.parquet} and
\texttt{results/<regime>\_gp.parquet}) after that regime's post-sweep QA
gate passed. The stable columns of every regime-indexed table are pending
placeholders---that regime's results files do not yet exist in final
form---and the fragments carry a first-line marker to that effect until
all three regimes are present. Each table's note states the regime, the
number of populations, and---where a contrast is involved---the number of
cells dropped by the common-cell restriction of Addendum~3 (\S11).}

\subsection{Validation gates}

Gates 1--3 (Section~\ref{sec:design-gates}) passed before any real-data
result was read. The one substantive finding from gate~1---the mild
over-dispersion of two-stage SVD Lee--Carter on Poisson data---is carried
into the reading of the SVD-LC rows below. A fourth, post-sweep gate
(\texttt{scripts/final\_qa.py}) classifies every cell that produced no
valid row before any table is generated; its output is the first table of
this section.

\subsection{What was scored: design-floor and method-failure cells}
\label{sec:results-infeasible}

% tab-infeasible is a page-spanning longtable fragment generated by
% scripts/make_tables.py; its caption and label live in the fragment.
\inputtable{tab-infeasible}

Table~\ref{tab:infeasible} lists every (population, family, mechanism) cell
of the primary grid that produced no valid row, in the classes of the
QA gate; the unabridged ledger, one row per affected cell with the full
recorded reason, is Appendix~\ref{app:infeasible}. The shift sweep
produced 2{,}000 cell rows, of which 210 are error rows: 126 structural,
56 method, 28 in the gate's residual class, and none machine---the gate
that refuses to generate tables while a broken-run row exists passed on
the final files.

\emph{Structural} cells are design-floor cells: the panel is shorter than
the mechanism's calibration or ensemble-member requirement, or falls under
the registered admissibility floor of $n_{\text{train}}\ge 15$ complete
training years (Addendum~3, \S4). The 126 structural rows trace the
training lengths of Table~\ref{tab:populations} exactly: KOR, the shortest
shift panel, appears under every mechanism with a calibration demand; HRV
joins it wherever the demand grows from the split wrapper's single
calibration block to the ten refits of EnbPI and the copula wrapper; CHL
enters under the largest windowed demands (the EnbPI and copula arms of
the sparse VAR, CNN-LC and LSTM), and the LSTM's construction alone adds
HKG.
These cells are a property of the design, not a failure of any method.

\emph{Method} cells are those in which a family's own sampler declines to
produce a predictive law on that panel under a registered rule. In the
shift regime this class is exclusively the sparse VAR's
explosive-coefficient-draw rejection (Addendum~3, \S7---rejection and
never clipping): it fires on 11 rows over six populations under the
native law (CHL, HKG, HRV, JPN, KOR, TWN), on 15 rows over nine
populations under the Poisson bootstrap, and on 8--11 rows in each
conformal arm; these rows are read below as a property of the family
under its own registered rule. The Poisson-composition overflow guard
never fired.

The 28 residual rows (class \emph{other}) all belong to the multi-output
GP: cells whose trailing block of complete years is shorter than the
family's minimum training block (Section~\ref{sec:design-families}; the
sixty-year cap of Addendum~4 is Section~\ref{sec:design-deviations}).
They are design-floor in character---a family-specific admissibility
requirement rather than a sampler failure---and are classed apart because
the floor is the family's own, not the registered mechanism floor.

Error rows of every class are excluded from every mean in this paper, and
every contrast between arms below is computed on the intersection of cells
in which each compared arm produced a valid row, with the intersection
size reported (Addendum~3, \S11).

\subsection{Calibration where nothing breaks: the stable regime}

The stable regime has not yet been run: the stable rows of every
regime-indexed table in this section are a pending placeholder until
\texttt{results/stable.parquet} exists, and no stable-derived quantity---
including the second, comparative clause of \Htwo\ and \Hfive---is stated
as a result in this draft.

\subsection{H1 as a harness-validity check: ranking by RMSE versus ranking by proper scores}

% tab-h1-rankings is a page-spanning longtable fragment generated by
% scripts/make_tables.py; its caption and label live in the fragment.
\inputtable{tab-h1-rankings}

Table~\ref{tab:h1-rankings} ranks the distributional arms---the
family--mechanism cells whose samples constitute a predictive
distribution---three times within the shift regime: by mean RMSE on log
rates, by mean CRPS on log rates and by mean Poisson log score. Two
exclusions are by construction. The conformal arms do not appear, because
their CRPS, log score and PIT are placeholder scores from
uniform-in-interval samples and are never ranked against a distributional
mechanism. CBD is ranked in a separate block: it is fit on ages 55--99
and therefore scored on 45 ages, and a mean over a different age support
is not the same quantity as a mean over ages 0--99. Means are taken
within population--sex first and then across populations, never
exposure-weighted, and the ranking is computed on the common cells of the
block (22 of 40 units kept; the dropped 18 trace the sparse VAR, GP and
LSTM failure cells of Table~\ref{tab:infeasible}).

\Hone\ as registered---rankings by RMSE and by CRPS disagree, with rank
correlation below one and at least one top-three inversion, in the shift
regime---is supported, in its mild form against CRPS and in a much
stronger form against the log score. The Spearman correlation between the
RMSE and CRPS orderings is $0.95$: below one, with the registered
top-three inversion present---the RMSE top three (LSTM-$\kappa_t$ under
dropout and ensemble, then SVD-LC native) share only one member with the
CRPS top three (the two sparse VAR arms, then LSTM dropout). Against the
Poisson log score the point-accuracy ordering nearly dissolves:
$\rho=0.12$. The arms with the best RMSE are among the worst by log
score---the LSTM ensemble is 2nd by RMSE and 16th of 18 by log score,
SVD-LC native 3rd and 12th---while the NB head, mid-pack by RMSE (ranks
10--12), takes the top three log-score positions together with CNN-LC
dropout, and neural-LC dropout, last by RMSE, is 5th by log score. A
point-accuracy league table and a density-forecast league table of the
same arms over the same break are close to unrelated, which is the
inversion the published literature reports
\citep{barigou2023bayesian,goes2025bayesian,schnurch2022point} recurring
under pre-registration.
% Intervals: results/h5_extras.json (scripts/h5_extras.py) -- cluster
% bootstrap resampling populations with replacement on the tab-h1
% common-cell block, B = 2000, seed 20260830.
Both correlations carry the cluster-bootstrap interval of
Section~\ref{sec:metrics-inference}, obtained by resampling with
replacement the thirteen populations that survive the common-cell
restriction and recomputing the block's means, ranks and correlations in
each of 2{,}000 replicates: $[0.77, 0.96]$ for $\rho(\text{RMSE},
\text{CRPS})$ and $[-0.16, 0.21]$ for $\rho(\text{RMSE}, \text{log
score})$. The first interval excludes one, so the registered
below-one clause does not depend on which populations entered the
block; the second contains zero, so at the resolution thirteen
correlated clusters afford, the log-score ordering is indistinguishable
from one unrelated to the RMSE ordering.

\subsection{H2: marginal coverage across the break}

\begin{table}[p]\centering
\caption{Empirical coverage of nominal 50/80/95\% central intervals and
mean 95\% interval score, by family, mechanism and regime, with the number
of cells behind each mean. Conformal mechanisms are scored at their
construction level only, so their 50\% and 80\% columns are blank by
design (Addendum~2, \S3). Error rows are excluded from every mean and
counted in Table~\ref{tab:infeasible}; a regime whose results file does
not yet exist appears as a pending row.}
\label{tab:h2-coverage}
\inputtable{tab-h2-coverage}
\end{table}

Table~\ref{tab:h2-coverage} reports, for every admissible family--mechanism
cell, the mean empirical coverage at 50, 80 and 95\%, the mean 95\%
interval score, and the number of cells behind each mean; conformal rows
carry the 95\% columns only, because a conformal wrapper constructs one
interval at one level and its coverage and interval score are read from
that interval's bounds. The interval score is the width-penalised
companion that keeps a vague interval from passing as a calibrated one.

\Htwo\ as registered has two clauses. The comparative clause---the
shortfall exceeds the stable-regime shortfall---awaits the stable regime
and is not tested here. The universal clause---nominal 95\% intervals
under-cover in the shift regime \emph{for every model family}---is
refuted as a universal, and the refutation is the more informative
result. The classical cores under their own predictive law do under-cover,
severely: Poisson Lee--Carter 0.624, APC/RH 0.654, SVD-LC 0.692, CBD
0.799 (on its 45-age support). But the distributional NB head covers
0.933 natively, 0.956 under the deep ensemble and 0.972 under dropout;
CNN-LC under dropout covers 0.962; the multi-output GP 0.905; the sparse
VAR 0.937--0.938 on the cells its rejection rule admits. Three of those
arms sit at or above nominal, so the direction of the miscalibration is
not even uniformly downward.

The registration deliberately declined a directional claim on
neural-versus-classical ordering, and the working conjecture that
motivated the audit---coverage degradation larger for the neural
families---is refuted in its strong form. The split does not run along
the architecture class at all: LSTM-$\kappa_t$ (0.647 ensemble, 0.685
dropout) and the neural-LC ensemble (0.694) break exactly like the
classical cores, while their sibling neural arms with a count likelihood
or dropout smoothing (NB head, CNN-LC dropout) are the best-covered arms
in the study. What separates the broken arms from the calibrated ones is
the family--mechanism pair, not the presence of a neural network---the
finding the crossed design exists to expose, taken up in
Section~\ref{sec:results-contrasts}.

The interval score orders the same failure: the broken natives pay twice,
in missed observations and in width---Poisson Lee--Carter native 4.456,
neural-LC ensemble 4.321, APC/RH 3.401, against NB native 1.684 and
CNN-LC dropout 1.677. The largest interval score in the table, neural-LC
dropout at 5.559 with coverage 0.811, is an arm buying part of its
coverage with indiscriminate width. Wrapping any family in a conformal
band moves its 95\% coverage into a narrow range regardless of how broken
the native law is: 0.848--0.944 for split and EnbPI, 0.939--0.981 for
the copula arm---the repair whose cost is priced in
Table~\ref{tab:dm-mcs}.

\subsection{H3: joint path coverage}

\begin{table}[p]\centering
\caption{Marginal coverage of nominal 95\% intervals, pooled over horizons,
against joint path coverage over the whole horizon set ($H=5$ in the shift
and stable regimes, $H=9$ in the placebo), and their difference, by family,
mechanism and regime. Conformal rows enter with both quantities read from
their interval bounds at the construction level (Addendum~3, \S6).}
\label{tab:h3-joint}
\inputtable{tab-h3-joint}
\end{table}

Table~\ref{tab:h3-joint} sets, for each family--mechanism cell, mean
marginal coverage at 95\% beside mean joint path coverage at the same
nominal level and the gap between the two: the difference between the
coverage a mortality paper reports and the coverage a liability depends
on. The registered direction of \Hthree---joint path coverage materially
below marginal for every family---holds in every row of the table: the
gap is negative for all 50 admissible arms, without exception. For the
broken classical natives the collapse is near-total: Poisson Lee--Carter
falls from 0.624 marginal to 0.298 joint, SVD-LC from 0.692 to 0.417,
APC/RH from 0.654 to 0.330, CBD from 0.799 to 0.504. The gap is not a
symptom of bad marginal calibration that good calibration removes: the
best-covered marginal arms lose it too---NB dropout falls from 0.972 to
0.897, CNN-LC dropout from 0.962 to 0.870---so an arm whose reported
marginal coverage is nominal can still miss between one in ten and one in
three whole trajectories.

The one mechanism constructed to be joint over $h=1,\dots,H$ behaves as
constructed: the copula-conformal arm holds joint path coverage between
0.805 (CBD) and 0.927 (SVD-LC) across the ten families, with gaps of
$-0.049$ to $-0.136$---the smallest in its family block in every family
but the NB head, whose dropout arm's gap of $-0.075$ is narrower still.
The registered comparator for \Hthree\ is not the marginal rate but each
model's own simulated-path-implied joint coverage (Addendum~3, \S8),
computed from the same predictive paths, so that the
realised-minus-implied gap can take either sign; that quantity is not
among the columns the runner emits. \todo{model-implied joint coverage
column (Addendum~3 \S8) is not emitted by the runner; add it to the
analysis stage or record the omission as a deviation in
Section~\ref{sec:design-deviations}}

\begin{figure}[htbp]\centering
\inputfigure{fig-joint-vs-marginal-shift}
\caption{Marginal 95\% coverage against joint path coverage over
$h=1,\dots,H$ as paired bars per mechanism, one panel per family, shift
regime, on the common cells of each panel (intersection size in the panel
title). Conformal arms are drawn in a distinct style and read from their
interval bounds; the nominal level is marked. The stable and placebo
counterparts are the same figure on those regimes' files.}
\label{fig:joint-vs-marginal}
\end{figure}

\subsection{H4: where in age the intervals fail}

\begin{table}[p]\centering
\caption{Empirical coverage of nominal 95\% intervals within the age bands
$0$--$24$, $25$--$64$ and $65$--$99$, by family, mechanism and regime.
CBD is fit on ages 55--99: its $0$--$24$ entry is blank and its
$25$--$64$ entry covers ages 55--64 only, as the table note states.
Conformal rows are read at the same bands the wrappers calibrate on.}
\label{tab:h4-age}
\inputtable{tab-h4-age}
\end{table}

Table~\ref{tab:h4-age} gives, for each family--mechanism cell, the 95\%
coverage within the three registered age bands, so that the $65$--$99$
and $25$--$64$ columns can be read against each other as \Hfour\
requires; the conformal wrappers are Mondrian-stratified on the same
bands, so a conformal row is read at the resolution it was calibrated at.
CBD contributes no $0$--$24$ value and a $25$--$64$ value restricted to
ages 55--64, which the table note records so that its row is not read
against the full-age families.

\Hfour\ as registered---shift-regime coverage at ages 65--99 lower than
at 25--64 \emph{for every family}---is refuted as a universal: the age
profile of the failure is family-specific, and no single clean pattern
spans the table. The registered direction does hold, sharply, for the
three families that forecast through an extrapolated period index and
broke in Table~\ref{tab:h2-coverage}: SVD-LC native falls from 0.725 at
25--64 to 0.518 at 65--99, APC/RH from 0.661 to 0.589, and
LSTM-$\kappa_t$ most steeply of all, 0.687 to 0.426 (ensemble) and 0.733
to 0.463 (dropout), against markedly better 0--24 rows in each case
(0.882--0.920 for SVD-LC and the LSTM, 0.732--0.744 for APC/RH): for
these arms the COVID miss is concentrated where COVID mortality was, at
the old ages. But Poisson Lee--Carter is uniformly
poor across all three bands (0.639 / 0.602 / 0.639)---its failure has no
age profile---and three arms reverse the registered direction outright:
CBD covers 0.849 at 65--99 against 0.626 at its 55--64 stub, CNN-LC
ensemble 0.880 against 0.847, and neural-LC dropout 0.934 against 0.874.
The neural-LC dropout arm instead collapses at the young ages: 0.538 at
0--24, the lowest band entry of any arm outside the index-driven
families. The calibrated families (NB head, sparse VAR, GP) are
near-flat in age, and the Mondrian-stratified wrappers flatten even the
broken families' profiles (SVD-LC split: 0.944 / 0.897 / 0.945). The
reversal anticipated by \citet{dowd2010backtesting}, which Addendum~3
(\S8) pre-committed to reporting as informative, is therefore present in
the table---but only for particular family--mechanism pairs, not as the
complement of a uniform old-age failure.

\begin{figure}[htbp]\centering
\inputfigure{fig-coverage-by-age-shift}
\caption{Single-age empirical coverage of nominal 95\% intervals (mean
over horizons), one panel per family, mechanisms as lines with conformal
arms dashed, shift regime, on the common cells of each panel. Ages a
family does not score are left blank and CBD's panel states its age
support; the horizontal line is the nominal level. The stable and placebo
counterparts are the same figure on those regimes' files.}
\label{fig:coverage-by-age}
\end{figure}

\subsection{PIT uniformity, reliability and the Murphy decomposition}

\begin{table}[htbp]\centering
\caption{Distributional arms only: mean Kolmogorov--Smirnov distance of the
randomised PIT from uniformity, and the share of population--sex cells
whose nominal KS $p$-value falls below $0.05$, by family, mechanism and
regime; conformal rows are left blank because their PIT is a placeholder.
The $p$-value is descriptive: PIT values within a cell are dependent
across ages and horizons (Addendum~2, \S4), and formal inference on
calibration is population-clustered.}
\label{tab:pit}
\inputtable{tab-pit}
\end{table}

\begin{table}[p]\centering
\caption{Murphy decomposition of the 95\% interval-hit Brier score
(reliability, resolution, uncertainty) and the PIT-scale reliability
against the uniform reference, by family, mechanism and regime; with one
constant nominal level the resolution term is zero by construction
(table footnote). PIT-scale column: distributional arms only.}
\label{tab:murphy}
\inputtable{tab-murphy}
\end{table}

Table~\ref{tab:pit} reports, for the distributional arms, the mean KS
distance of the randomised PIT from the uniform and the share of cells
whose nominal KS $p$-value is below $0.05$; the $p$-value is descriptive
only, because PIT values across ages and horizons within a cell are
dependent, and the shape of each departure is read from
Figure~\ref{fig:pit-hist}. The ordering matches
Table~\ref{tab:h2-coverage}---the smallest KS distances belong to the GP
(0.172), the NB head (0.182--0.189) and the sparse VAR (0.188), the
largest to Poisson Lee--Carter (0.374) and APC/RH (0.325)---but the
levels carry the sharper message: no arm is distributionally calibrated
through the break. Even for the best arms the share of cells whose
nominal KS $p$-value falls below 0.05 never drops beneath 0.844, and for
most arms it is 0.925--1.000. Coverage at a single level is a much
weaker property than PIT uniformity: the NB head's 95\% interval can be
essentially nominal while its predictive distribution, read over all
levels at once, is still rejected as uniform in nearly every cell.

Table~\ref{tab:murphy} decomposes the Brier score of the 95\% hit
indicator: with a constant forecast probability there is one bin, so
reliability is the squared coverage gap, resolution is zero by
construction, and uncertainty is the base-rate term. The reliability
column spans two orders of magnitude across arms---0.1515 for Poisson
Lee--Carter native, 0.1311 for APC/RH and 0.1226 for the LSTM ensemble
at the broken end, against 0.0015 for CNN-LC dropout, 0.0022 for NB
dropout and 0.0012--0.0223 for the conformal rows---which restates the
coverage table as a proper-score share. The PIT-scale reliability beside
it, the chi-square-type distance from uniformity that a single level's
coverage cannot see, orders the arms the same way (Poisson Lee--Carter
0.1904 down to NB dropout 0.0302) but never approaches zero, in
agreement with the KS column of Table~\ref{tab:pit}: the calibrated arms
are calibrated at 95\%, not everywhere. \todo{Section~\ref{sec:metrics}
describes a Murphy decomposition of the interval score and of exceedance
Brier scores with an isotonic recalibration; the harness emits the
one-bin hit-indicator form and the PIT-scale form above---reconcile the
metrics text with the emitted columns before submission}

\begin{figure}[htbp]\centering
\inputfigure{fig-pit-hist-shift}
\caption{Ten-bin randomised PIT histograms averaged across population--sex
cells, in a family $\times$ distributional-mechanism grid, shift regime;
the horizontal line is the uniform reference. Conformal arms are absent
because their PIT is a placeholder. A U shape reads as over-confidence, a
hump as under-confidence, a slope as bias. The stable and placebo
counterparts are the same figure on those regimes' files.}
\label{fig:pit-hist}
\end{figure}

\begin{figure}[htbp]\centering
\inputfigure{fig-reliability-shift}
\caption{Reliability diagram: empirical coverage against nominal level
(50, 80, 95\%) per mechanism, one panel per family, shift regime, with
the diagonal as perfect calibration. Conformal arms contribute a single
point at their construction level, 95\%, in a distinct style. The stable
and placebo counterparts are the same figure on those regimes' files.}
\label{fig:reliability}
\end{figure}

\subsection{Forecast comparison: Diebold--Mariano and the model confidence set}

% tab-dm-mcs is a page-spanning longtable fragment generated by
% scripts/make_tables.py; it carries its own \caption and \label{tab:dm-mcs}.
\inputtable{tab-dm-mcs}

Table~\ref{tab:dm-mcs} shows the two registered inference procedures, and
its loss column is the reason it is one table: a contrast among
distributional arms is decided on CRPS on log rates, while any contrast
that includes a conformal arm is decided on the per-horizon 95\% interval
score, because a conformal cell's CRPS is a placeholder. Every entry is
computed on the intersection of cells in which each compared arm produced
a valid row, and the kept and dropped counts are part of the table.

The model confidence set over the classical families under their own
predictive law is formed over the four full-age families only: the
five-family set including CBD is skipped, with its reason recorded in the
table, because a mean CRPS over 45 ages is not comparable with one over
100 and the registration forbids working around an incomparable loss. On
the 29 common cells the 90\% set retains \{SVD-LC native, sparse VAR
native\} (elimination $p$-values 0.13 and 1.00), eliminating Poisson
Lee--Carter and APC/RH at $p=0.00$---with the caveat, carried from
Table~\ref{tab:infeasible}, that the sparse VAR is scored only on the
cells its own rejection rule admits.

The within-family conformal sets on the interval score split three ways.
In five families---SVD-LC, LSTM-$\kappa_t$, NB head, neural-LC and sparse
VAR---EnbPI is the sole survivor, eliminating the split and copula arms
at $p$-values between 0.00 and 0.08; in Poisson Lee--Carter the copula
arm is the sole survivor ($p=0.01$ against both others); and in CBD,
CNN-LC, the GP and APC/RH all three wrappers are retained. Read with
Table~\ref{tab:h3-joint}, the EnbPI wins are a width result, not a
calibration result: the interval score at a single level rewards the
narrowest band that roughly covers, which EnbPI's residual recycling
produces, while the copula arm's wider joint band is what actually holds
path coverage. The two registered procedures answer different questions
and disagree accordingly.

The Diebold--Mariano block prices the conformal repair of
Table~\ref{tab:h2-coverage} and is the attribution finding of the crossed
design. The contrast is defined for the seven families with a native
arm. Split conformal is significantly \emph{better} than the native law
for exactly the four classical cores whose native intervals broke:
Poisson Lee--Carter ($\Delta=1.895$, $p=0.004$), APC/RH ($\Delta=1.039$,
$p=0.013$), SVD-LC ($\Delta=0.700$, $p=0.018$) and CBD ($\Delta=0.318$,
$p=0.017$), with $G=19$ population clusters each. It is significantly
\emph{worse} for exactly the three families whose native law was already
calibrated: the NB head ($\Delta=-1.024$, $p=0.000$, $G=19$), the
multi-output GP ($\Delta=-0.413$, $p=0.009$, $G=16$) and the sparse VAR
($\Delta=-0.318$, $p=0.000$, $G=14$). No family is left where the
wrapper is free: conformalisation is insurance that pays out precisely
when the model's own uncertainty was wrong and costs a premium precisely
when it was right. Which side of the line a family falls on is a
property of the family--mechanism pair, not of the wrapper---the
sharpest form of the study's central claim that architecture and
uncertainty mechanism are separate factors.

\subsection{Twin crises: 1914--1922 beside 2020--2024}

\begin{table}[htbp]\centering
\caption{Placebo (train to 1913, test 1914--1922, eleven populations) and
shift (train to 2019, test 2020--2024, twenty populations) side by side:
marginal 95\% and joint path coverage by family and mechanism, with the
placebo strata of Addendum~1 as a further column. Joint coverage spans
$H=9$ placebo and $H=5$ shift horizons---not one scale. Sparse VAR's
empty placebo cells are a finding, not a gap: every native draw on the
pre-1914 panels tripped the registered explosive-path rejection
(Table~\ref{tab:infeasible}).}
\label{tab:twin-crises}
\inputtable{tab-twin-crises}
\end{table}

Table~\ref{tab:twin-crises} sets the 1914--1922 placebo beside the
2020--2024 shift. Read within regime (the joint columns span $H=9$
against $H=5$), the placebo reproduces the shift regime's family split a
century earlier: the full-age classical cores under-cover (Lee--Carter
0.787, Poisson Lee--Carter 0.780, APC/RH 0.761 marginal against a
nominal 0.95), the point-forecast neural mechanisms are worst
(neural-LC ensemble 0.661, LSTM 0.636), and the families that came
through COVID calibrated sit near the top of the placebo ordering too
(NB head 0.849, multi-output GP 0.820). Joint path coverage collapses
in the same order---0.146 (neural-LC ensemble), 0.172 (LSTM),
0.293--0.335 (the classical cores) over nine horizons---so the H3 gap
is not a COVID artefact either. Sparse VAR is the one family with no
placebo verdict: every native draw on the pre-1914 panels tripped the
registered explosive-path rejection (zero valid rows,
Table~\ref{tab:infeasible}), which also empties the common-cell
intersection (addendum 3~\S11) of the registered classical model
confidence set in this regime.

The exception is CBD, and it is the most instructive cell in the table:
near-nominal through the placebo break (0.986 marginal, 0.918 joint)
and broken through the COVID one (0.799, 0.504). The 1918 influenza
concentrated its excess mortality at young-adult ages, below CBD's
age-55 support, while COVID's excess fell at the old ages CBD models.
Whether an interval survives a crisis therefore depends not only on the
family's dynamics but on where in the age range the crisis lands---the
age-localisation of miscalibration that H4 establishes \emph{within}
the shift regime, operating here \emph{between} regimes.

The Addendum~1 strata sharpen the reading. In the neutral
register-based populations the classical cores hold 0.86--0.89---the
1918 flu alone bends them without breaking them---while the
belligerent-total stratum, where war losses compound the flu, drops the
same families to 0.55--0.59, with the civilian-only stratum in between.
Within one family and one regime, coverage failure scales with the size
of the break.

The paired mechanism contrasts, however, do not transfer. The
split-conformal rescue that is decisive under COVID
(Section~\ref{sec:results-contrasts}) is absent here: on the interval
score, native versus split is insignificant for every full-age
classical core (Lee--Carter $\Delta=-0.091$, $p=0.689$; Poisson
Lee--Carter $p=0.127$; APC/RH $p=0.563$), and the NB head's native law
significantly beats its own wrapper ($\Delta=-0.461$, $p=0.008$)---a
calibration window drawn from the pre-1914 years knows as little about
the coming shock as the model it corrects. The ensemble-versus-dropout
verdicts flip family by family between the two crises: CNN-LC,
indistinguishable under COVID ($p=0.739$), favours dropout here
($\Delta=+0.531$, $p=0.0002$); LSTM, which favours dropout under COVID
($p=0.0006$), is indistinguishable here ($p=0.478$). A family's
fragility is predictable from its past crisis; which uncertainty
mechanism patches it best is not.

\subsection{Architecture or mechanism? Contrasts within sub-grids}
\label{sec:results-contrasts}

The payoff of the crossed design. Each contrast below is pre-registered and
computed as a paired difference within the admissible sub-grid named, on
the common cells of that sub-grid.

\begin{enumerate}
\item \emph{Conformal versus native.} Does wrapping a fit in a conformal
  band change its shift-regime interval score and coverage? The native
  versus split-conformal block of Table~\ref{tab:dm-mcs} carries this
  contrast, per family, on the interval score, and its answer is signed by
  the family: the wrapper significantly improves the four classical cores
  whose native law broke and significantly worsens the NB head, GP and
  sparse VAR, whose native law had not. The width side of the question is
  read from the interval-score column of Table~\ref{tab:h2-coverage}: for
  the broken families the wrapper buys its coverage while \emph{shrinking}
  the score (Poisson Lee--Carter 4.456 native against 2.674 split), for
  the calibrated ones it pays in width (NB head 1.684 native against
  2.717 split).
\item \emph{Ensemble versus Monte Carlo dropout, neural families only.} Do
  two mechanisms on the same architecture give the same coverage?
  Descriptively, Table~\ref{tab:h2-coverage} answers no---on the same
  CNN-LC architecture, dropout covers 0.962 against the ensemble's 0.882,
  and on neural-LC 0.811 against 0.694. The registered paired test agrees
  and adds a sign that the descriptive reading does not supply: on the
  interval score the two mechanisms differ significantly on two of the four
  neural families, \emph{in opposite directions}. Dropout is favoured on
  LSTM-on-$\kappa_t$ ($\Delta = +0.224$, $t = 3.78$, $p = 0.0006$) and the
  ensemble on neural-LC ($\Delta = -1.238$, $t = -4.18$, $p < 0.0001$),
  while CNN-LC ($p = 0.739$) and the NB head ($p = 0.139$) are
  indistinguishable. The mechanism is therefore not separable from the
  architecture it wraps: no single ordering of the two survives a change of
  family, which is the crossed design's reason for existing.
\item \emph{Bootstrap versus native, classical families only.} Does adding
  parameter uncertainty move a classical interval toward nominal at the
  break? Descriptively the movement is positive and small (SVD-LC 0.692
  to 0.727, CBD 0.799 to 0.830 in Table~\ref{tab:h2-coverage}), nowhere
  near the repair a conformal wrapper delivers. The registered paired test
  makes the smallness precise: the semiparametric bootstrap improves the
  interval score significantly on three of the five classical families---
  SVD-LC ($\Delta = -0.112$, $p < 0.0001$), Poisson Lee--Carter
  ($\Delta = -0.115$, $p < 0.0001$) and Renshaw--Haberman
  ($\Delta = -0.101$, $p = 0.0002$)---leaves CBD unchanged ($p = 0.681$),
  and significantly \emph{worsens} sparse VAR ($\Delta = +0.168$,
  $p = 0.0008$, on the 22 cells where its native arm survives the
  explosive-draw rejection of Addendum~3~\S7). The three improvements are
  an order of magnitude smaller than the conformal repair of the same
  families in the contrast above (for Poisson Lee--Carter, $-0.115$ against
  $+1.895$ in favour of the wrapper): resampling the estimator moves a
  broken interval, but it does not fix it.
\item \emph{Joint-path conformal versus split conformal, all families.} The
  \Hthree\ mechanism contrast: the within-family conformal confidence sets
  of Table~\ref{tab:dm-mcs} on the interval score, read with the joint
  column of Table~\ref{tab:h3-joint}. The two tables pull in opposite
  directions by construction: the copula arm holds joint path coverage of
  0.805--0.927 where the split arm holds 0.692--0.862, yet on the
  single-level interval score the copula arm survives as sole set member
  only for Poisson Lee--Carter and is eliminated in five families where
  EnbPI's narrower band wins. A practitioner who needs the path pays for
  it at the level; the study reports the price rather than adjudicating
  it.
\end{enumerate}
